% Source: https://artofproblemsolving.com/community/c3h2751389_800th_post_1_year_aops_anniversary_and_maff_problems
% Stable identity: aops-topic-2751389-q02
The six different values form a divisibility chain $l=d_1<d_2<\dots<d_6=g$ with $d_i\mid d_{i+1}$. Let $v$ be the repeated value; the ten numbers are $d_1,\dots,d_6$ together with $v$ four more times. The condition becomes
\[g-l=4v+d_2+d_3+d_4+d_5.\]
If $v\ne l$, then $v$ and $d_2,\dots,d_6$ are all multiples of $d_2$, and the equation $l=g-4v-(d_2+\dots+d_5)$ makes $l$ a multiple of $d_2>l$, which is impossible. So $v=l$ and
\[g=5l+d_2+d_3+d_4+d_5.\]
Each $d_i$ is a proper multiple of the previous one. Reducing modulo $d_2$: $d_2\mid5l$, and since $d_2=rl$ with $r\ge2$, $r\mid5$, so $d_2=5l$. Modulo $d_3$: $d_3\mid5l+d_2=10l$, and $d_3$ is a proper multiple of $5l$, so $d_3=10l$. Likewise $d_4\mid20l$ gives $d_4=20l$, and $d_5\mid40l$ gives $d_5=40l$. Then
\[g=5l+5l+10l+20l+40l=80l,\]
so $\dfrac gl=\boxed{80}$. (Check: with $l=1$, the numbers $1,1,1,1,1,5,10,20,40,80$ sum to $160$, and $160-81=79=g-l$.)
